\begin{table}[t]\centering\small
\begin{tabular}{lllrrr}\toprule
Model & Condition & Access & Detected & LPIPS & Class CLIP\\\midrule
RAR & Marked original & -- & 100.0\% & 0.0000 & 0.3030 \\
RAR & Clean LoRA 2,500 & W & 4.9\% & 0.6692 & 0.3044 \\
RAR & Clean full 2,500 & W & 7.2\% & 0.6618 & 0.3032 \\
RAR & NFPA 40 & O & 21.5\% & 0.4221 & 0.2630 \\
DCAE & Marked original & -- & 100.0\% & 0.0000 & 0.2934 \\
DCAE & Clean LoRA 2,500 & W & 0.2\% & 0.1510 & 0.2840 \\
DCAE & Clean full 2,500 & W & 0.4\% & 0.2148 & 0.2816 \\
DCAE & NFPA 40 & O & 33.6\% & 0.5082 & 0.2664 \\
\bottomrule\end{tabular}
\caption{DistSeal image-watermark removal on RAR-XL at 256 pixels and UViT-H/DC-AE at 512 pixels, with the original clean decoder and 512 matched samples per condition. Quality means come from the later matched native-resolution audit, consistently across weight and output interventions. Detected is the percentage retaining owner detection after the named intervention. Strict bit-accuracy thresholds are $>0.640625$ for RAR and $>0.6875$ for DC-AE. Clean LoRA and full training use 2,500 clean-data updates; NFPA strength 40 transforms the source output using independent SD2.1. Access: W requires the marked generator weights and clean training data; O requires only saved target outputs; -- marks the baseline. LPIPS is mean candidate-to-source perceptual distance, with lower values indicating closer pairs. Class CLIP is mean candidate alignment to the requested ImageNet class using CLIP ViT-B/32; it measures class agreement and does not certify retention of the original scene.}\end{table}
